%! Adding \& Subtracting Rational Expressions — Unit 5, Lesson 5.3 — Group Activity
\documentclass[10pt]{article}
\usepackage{algebra2-article}
\usepackage{algebra2-boxes}
\newcommand{\TallMath}[1]{$\displaystyle #1\rule[-1.4em]{0pt}{3.2em}$}

\begin{document}

\pageheader{Unit 5, Lesson 5.3}{Group Activity}
\namepartnerperiod

\vspace{0.08in}

\begin{headlinebox}{forestbg}
\small\textbf{Build the Denominator, Guard the Sign.} Two habits carry everything today. Factor every
denominator \emph{before} you decide what the LCD is --- it is built out of factors, never out of
products. And every time you subtract, put the second numerator in parentheses \emph{first}. As always,
an answer without its restrictions is half an answer.
\end{headlinebox}

\vspace{0.06in}

\begin{tcolorbox}[colback=white, colframe=black!40, title=\textbf{Tier R --- Remediate},
    fonttitle=\bfseries, arc=1.5mm, left=3mm, right=3mm, top=2mm, bottom=2mm, breakable]
Combine and simplify. State every restriction.

\vspace{4pt}
\begin{enumerate}[leftmargin=1.9em, itemsep=9pt]
  \item $\dfrac{2x}{x-6}+\dfrac{5}{x-6}$ \; (denominators already match) \\[3pt]
        Simplest form: \blank{2.2cm} \quad $x\neq$ \blank{1.2cm}
  \item $\dfrac{4x+1}{x+2}-\dfrac{x-8}{x+2}$ \; Write the numerator with parentheses first:
        $(4x+1)-(\rule{1.8cm}{0.4pt})$ \\[3pt]
        Numerator, simplified and factored: \blank{2.6cm} \quad Simplest form: \blank{2.4cm}
        \quad $x\neq$ \blank{1.2cm}
  \item $\dfrac{3}{x}+\dfrac{5}{2x}$ \\[3pt]
        LCD: \blank{1.4cm} \quad Building factor for the first fraction: \blank{1.2cm} \quad
        Answer: \blank{1.8cm}, \; $x\neq$ \blank{1.0cm}
  \item $\dfrac{1}{x-1}+\dfrac{2}{x+3}$ \\[3pt]
        LCD: \blank{2.6cm} \quad Answer: \blank{3.0cm} \quad $x\neq$ \blank{1.8cm}
  \item In problem 2, do it the wrong way on purpose: drop the parentheses and write
        $4x+1-x-8$. \\[3pt]
        Wrong numerator: \blank{1.8cm} \quad Now test both at $x=0$: the original expression equals
        \blank{1.2cm}, your correct answer equals \blank{1.2cm}, and the wrong one equals
        \blank{1.2cm}.
\end{enumerate}
\end{tcolorbox}

\begin{tcolorbox}[colback=white, colframe=black!40, title=\textbf{Tier A --- Approaching Proficiency},
    fonttitle=\bfseries, arc=1.5mm, left=3mm, right=3mm, top=2mm, bottom=2mm, breakable]
\textbf{Part 1 --- The full table.} Fill in every cell. One row hides a pair of \emph{opposite}
binomials, and one row simplifies further than you expect.
\begin{center}
\renewcommand{\arraystretch}{1.7}
\begin{tabularx}{\linewidth}{@{}l X X X@{}}
\hline
\textbf{Expression} & \textbf{Denominators, factored} & \textbf{LCD \& restrictions} & \textbf{Simplest form} \\
\hline
\TallMath{\frac{7}{4x^2}-\frac{3}{10x}} & \blank{2.2cm} & \blank{2.2cm} & \blank{2.2cm} \\
\TallMath{\frac{x}{x^2-4}+\frac{1}{x+2}} & \blank{2.2cm} & \blank{2.2cm} & \blank{2.2cm} \\
\TallMath{\frac{2}{x-4}-\frac{16}{x^2-16}} & \blank{2.2cm} & \blank{2.2cm} & \blank{2.2cm} \\
\TallMath{\frac{5}{x-2}+\frac{3}{2-x}} & \blank{2.2cm} & \blank{2.2cm} & \blank{2.2cm} \\
\hline
\end{tabularx}
\end{center}

\vspace{4pt}
\textbf{Part 2 --- Error analysis.} A student turns in this work:
\begin{center}
\fbox{\parbox{0.66\linewidth}{\centering \vspace{2pt}
$\dfrac{1}{x} + \dfrac{1}{3} \;=\; \dfrac{1+1}{x+3} \;=\; \dfrac{2}{x+3}$ \\[3pt]
{\small ``I added the tops and I added the bottoms.''}\vspace{2pt}}}
\end{center}
\begin{enumerate}[label=(\alph*), leftmargin=1.9em, itemsep=6pt]
  \item Disprove it with a number. At $x=1$: the original equals \blank{1.4cm} and the student's
        answer equals \blank{1.4cm}.
  \item Do it correctly: LCD $=$ \blank{1.4cm}, so $\dfrac1x+\dfrac13=$ \blank{2.2cm}, \;
        $x\neq$ \blank{1.0cm}.
  \item Multiplying rational expressions \emph{does} let you work straight across the tops and the
        bottoms. Adding does not. In one sentence, what is different about adding?
        \writelines{2}
\end{enumerate}
\end{tcolorbox}

\begin{tcolorbox}[colback=white, colframe=black!40, title=\textbf{Tier E --- Extension},
    fonttitle=\bfseries, arc=1.5mm, left=3mm, right=3mm, top=2mm, bottom=2mm, breakable]
\textbf{Part 1 --- One complex fraction, two roads.} Simplify
\[
  \frac{\;\dfrac{1}{x}-\dfrac{1}{4}\;}{\;\dfrac{1}{x^2}-\dfrac{1}{16}\;}
\]
\begin{enumerate}[label=(\alph*), leftmargin=1.9em, itemsep=6pt]
  \item \textbf{Method 1 (combine, then divide).} Top: \blank{2.0cm} \quad Bottom: \blank{2.6cm}
        \\[3pt]
        Keep--change--flip and simplify: \blank{2.4cm}
  \item \textbf{Method 2 (multiply through by the LCD of the little fractions).} That LCD is
        \blank{1.4cm}. \\[3pt]
        Top becomes \blank{2.2cm}; bottom becomes \blank{2.2cm}; answer: \blank{2.0cm}
  \item Complete restriction list: $x\neq$ \blank{2.4cm}. Which value(s) came from the \emph{main
        fraction bar}? \blank{2.0cm}
  \item One method leaves \emph{every} restriction readable somewhere on the page; the other erases
        one of them. Which method erases which restriction, and what step destroyed the evidence?
        \writelines{2}
\end{enumerate}

\vspace{5pt}
\textbf{Part 2 --- Build one backwards.} Write a \emph{difference} of two rational expressions whose
simplified value is $\dfrac{2}{x+1}$ and whose restrictions are exactly $x\neq-1$ and $x\neq4$.
\begin{center}
\blank{3.6cm} $-$ \blank{3.6cm}
\end{center}
\begin{enumerate}[label=(\alph*), leftmargin=1.9em, itemsep=5pt]
  \item What is the LCD of your two denominators, and why did it have to contain $(x-4)$?
  \item On the graph of your expression, which of the two excluded values is a \textbf{hole} and which
        is a \textbf{wall}? How do you know without graphing anything?
\end{enumerate}
\writelines{2}

\vspace{5pt}
\textbf{Part 3 --- Two crews, one mural.} Crew A can paint a mural alone in $x$ hours. Crew B is slower
and needs $x+3$ hours.
\begin{enumerate}[label=(\alph*), leftmargin=1.9em, itemsep=6pt]
  \item In one hour, Crew A finishes \blank{1.0cm} of the mural and Crew B finishes \blank{1.4cm} of
        it. Working together, they finish \blank{2.6cm} of it per hour.
  \item Time is $1$ whole mural divided by their combined rate --- a \textbf{complex fraction}.
        Simplify it: \\[3pt]
        Time together $=$ \blank{3.0cm} hours.
  \item If Crew A alone takes $x=6$ hours, the two crews together take \blank{1.6cm} hours. Is that
        answer smaller than both $6$ and $9$? \blank{1.0cm} Why must it be? \blank{4.0cm}
  \item The algebra excludes $x=0$, $x=-3$, and $x=-\tfrac32$. Which values of $x$ actually make sense
        for two paint crews? \blank{3.0cm}
\end{enumerate}
\end{tcolorbox}

\end{document}
